This chapter defines the architectural meaning of each instruction. Opcode
assignments and binary layouts are defined separately by the instruction
encoding and opcode map.

\section{Notation}

\subsection{Instruction Notation}

\begin{itemize}
\item \code{\$}: prefix for general-purpose registers.
\item \code{\#}: prefix for floating-point registers.
\item \code{\%}: prefix for named control registers.
\item \code{rd}: destination register.
\item \code{rd0}: first destination register.
\item \code{rd1}: second destination register.
\item \code{ra}: first source register.
\item \code{rb}: second source register.
\item \code{rs}: source register.
\item \code{immN}: immediate N-bit value encoded in the instruction.
\item \code{offsetN}: signed N-bit \code{\%pc}- or base-relative offset.
\item \code{x.y}: bit \code{y} of register \code{x}.
\item \code{\%pc}: program counter.
\item \code{\%sr}: processor status register.
\item \code{N}: negative flag.
\item \code{Z}: zero flag.
\item \code{C}: carry flag.
\item \code{V}: overflow flag.
\item \code{I}: infinity flag.
\item \code{U}: underflow flag.
\item \code{X}: inexact flag.
\item \code{E}: arithmetic error flag.
\end{itemize}

\subsection{Assignment Notation}

\begin{itemize}
\item The notation \code{\$rd}~\assign~\code{x} means that value \code{x} is
written to register \code{\$rd}.
\item The notation \code{\%pc}~\assign~\code{x} means that execution continues
from address \code{x}.
\item The notation \code{[x]} means memory at address \code{x}.
\item The notation \code{SEXT(x)} means sign-extension of \code{x} to 32 bits.
\item The notation \code{ZEXT(x)} means zero-extension of \code{x} to 32 bits.
\item The notation \code{a:b} means a register pair, with \code{a} holding the
upper 32 bits and \code{b} holding the lower 32 bits unless otherwise specified.
\end{itemize}

\subsection{Immediate Values}

\begin{itemize}
\item Unless otherwise specified, immediate operands are sign-extended to 32
bits before use.
\item Arithmetic immediate instructions sign-extend their immediate operands.
\item Logical immediate instructions zero-extend their immediate operands.
\item Shift counts are interpreted as unsigned values.
\item Bit indexes are interpreted as unsigned values in the range 0--31 for
general purpose registers, and 0--63 for floating-point registers.
\end{itemize}

\subsection{Conditional Suffixes}

\begin{itemize}
\item Conditional instruction mnemonics use a dot-separated condition suffix.
\item Register-comparison branches use the form \code{b[d].<cond>}, where
\code{d} signifies the operands are floating-point registers.
\item Flag-based branches use the form \code{bf.<cond>}.
\item Future conditional instructions should use the same suffix convention.
\end{itemize}

\begin{center}
\begin{tabular}{lll}
\hline
Suffix & Meaning & Basis \\
\hline
\code{.eq}  & equal & register comparison \\
\code{.ne}  & not equal & register comparison \\
\code{.lt}  & signed less than & signed register comparison \\
\code{.le}  & signed less than or equal & signed register comparison \\
\code{.gt}  & signed greater than & signed register comparison \\
\code{.ge}  & signed greater than or equal & signed register comparison \\
\code{.ltu} & unsigned less than & unsigned register comparison \\
\code{.leu} & unsigned less than or equal & unsigned register comparison \\
\code{.gtu} & unsigned greater than & unsigned register comparison \\
\code{.geu} & unsigned greater than or equal & unsigned register comparison \\
\hline
\end{tabular}
\end{center}

Flag-based branches use a separate set of suffixes:

\begin{center}
\begin{tabular}{lll}
\hline
Suffix & Meaning & Basis \\
\hline
\code{.ns} & negative set & \code{N~=~1} \\
\code{.nc} & negative clear & \code{N~=~0} \\
\code{.zs} & zero set & \code{Z~=~1} \\
\code{.zc} & zero clear & \code{Z~=~0} \\
\code{.cs} & carry set & \code{C~=~1} \\
\code{.cc} & carry clear & \code{C~=~0} \\
\code{.vs} & overflow set & \code{V~=~1} \\
\code{.vc} & overflow clear & \code{V~=~0} \\
\code{.es} & arithmetic error set & \code{E~=~1} \\
\code{.ec} & arithmetic error clear & \code{E~=~0} \\
\code{.is} & infinity set & \code{I~=~1} \\
\code{.ic} & infinity clear & \code{I~=~0} \\
\code{.us} & underflow set & \code{U~=~1} \\
\code{.uc} & underflow clear & \code{U~=~0} \\
\code{.xs} & inexact set & \code{X~=~1} \\
\code{.xc} & inexact clear & \code{X~=~0} \\
\hline
\end{tabular}
\end{center}

\section{Flag Behaviour}

Condition flags describe the result or status of the most recent flag-setting
operation.

\begin{itemize}
\item Arithmetic and logical instructions update condition flags by default.
\item Comparison instructions update condition flags without storing a result.
\item Load instructions update condition flags according to the resulting
register value.
\item Store instructions do not update condition flags.
\item Branch, jump, call, and return instructions do not update condition flags.
\item System instructions update condition flags only when explicitly specified.
\end{itemize}

Unless otherwise specified by an instruction, all condition flags are updated
or cleared according to the rules in this section.

\subsection{Negative Flag}

The negative flag \code{N} indicates that the numeric result is negative.

\begin{itemize}
\item For integer results, \code{N} is set when bit 31 of the result is set.
\item For floating-point results, \code{N} is set when the numeric binary64
result is negative.
\item Floating-point \code{+0.0} and \code{-0.0} do not set \code{N}.
\end{itemize}

\subsection{Zero Flag}

The zero flag \code{Z} indicates that the numeric result is zero.

\begin{itemize}
\item For integer and logical results, \code{Z} is set when all result bits are
zero.
\item For floating-point results, \code{Z} is set when the result is either
\code{+0.0} or \code{-0.0}.
\item Otherwise, \code{Z} is cleared.
\end{itemize}

\subsection{Carry Flag}

The carry flag \code{C} is used by integer arithmetic.

\begin{itemize}
\item For integer addition, \code{C} is set when an unsigned carry is produced
from bit 31.
\item For integer subtraction, \code{C} is set when no unsigned borrow is
required and cleared when a borrow is required.
\item For successful integer division, \code{C} is set when the remainder is
non-zero.
\item Floating-point operations clear \code{C} unless otherwise specified.
\item Logical operations clear \code{C} unless otherwise specified.
\end{itemize}

\subsection{Overflow Flag}

The overflow flag \code{V} is used by integer arithmetic.

\begin{itemize}
\item For signed integer arithmetic, \code{V} is set when the mathematical
result cannot be represented as a signed 32-bit two's-complement integer.
\item Floating-point operations clear \code{V} unless otherwise specified.
\item Logical operations clear \code{V} unless otherwise specified.
\end{itemize}

\subsection{Arithmetic Error Flag}

The arithmetic error flag \code{E} indicates an exceptional or invalid
arithmetic or conversion condition that does not trap in VIV-32 v\version{}.

\begin{itemize}
\item Integer divide-by-zero sets \code{E}.
\item Integer division that cannot produce a representable quotient sets
\code{E}.
\item Floating-point divide-by-zero sets \code{E}.
\item A floating-point operation producing \code{NaN} sets \code{E}.
\item An unordered floating-point comparison involving \code{NaN} sets
\code{E}.
\item A floating-point to integer conversion of \code{NaN}, infinity, or a
value outside the destination integer range sets \code{E}.
\item Arithmetic and conversion operations producing an ordinary valid result
clear \code{E}, unless otherwise specified.
\end{itemize}

\subsection{Infinity Flag}

The infinity flag \code{I} indicates that a floating-point result is infinite.

\begin{itemize}
\item \code{I} is set when the resulting floating-point value is positive or
negative infinity.
\item The presence of infinity does not by itself imply an arithmetic error;
\code{E} is determined by the operation that produced or processed the
value.
\item Integer operations clear \code{I} unless otherwise specified.
\end{itemize}

\subsection{Underflow Flag}

The underflow flag \code{U} indicates that a floating-point result is too small
in magnitude to be represented normally in the destination floating-point
format.

\begin{itemize}
\item \code{U} is set when a floating-point arithmetic or conversion result
enters the underflow range of its destination format.
\item For binary64 arithmetic, the destination format is binary64.
\item For binary64-to-binary32 conversion, the destination format is binary32.
\item Integer operations clear \code{U} unless otherwise specified.
\end{itemize}

\subsection{Inexact Flag}

The inexact flag \code{X} indicates that a numeric result required rounding.

\begin{itemize}
\item \code{X} is set when the mathematically exact floating-point or conversion
result cannot be represented exactly in the destination format.
\item \code{X} may be set independently of \code{U}, \code{I}, or \code{E}.
\item Integer operations clear \code{X} unless otherwise specified.
\end{itemize}

\section{Arithmetic Instructions}

Arithmetic operations can be performed on both general purpose and
floating-point registers. General purpose, single-target-register operations
wrap modulo \(2^{32}\).

\subsection{\code{add, addd}}

\begin{itemize}
\item Syntax: \code{add \$rd, \$ra, \$rb} | \code{addd \#rd, \#ra, \#rb}
\item Operation: \code{rd~\assign~ra~+~rb}
\end{itemize}

\subsection{\code{addi}}

\begin{itemize}
\item Syntax: \code{addi \$rd, \$ra, imm16}
\item Operation: \code{rd~\assign~ra~+~SEXT(imm16)}
\end{itemize}

\subsection{\code{sub, subd}}

\begin{itemize}
\item Syntax: \code{sub \$rd, \$ra, \$rb | subd \#rd, \#ra, \#rb}
\item Operation: \code{rd~\assign~ra~--~rb}
\end{itemize}

\subsection{\code{subi}}

\begin{itemize}
\item Syntax: \code{subi \$rd, \$ra, \$imm16}
\item Operation: \code{rd~\assign~ra~--~SEXT(imm16)}
\end{itemize}

\subsection{\code{neg, negd}}

\begin{itemize}
\item Syntax: \code{neg \$rd, \$ra | negd \#rd, \#ra}
\item Operation: \code{rd~\assign~0~--~ra}
\end{itemize}

\section{Logical Instructions}

Logical instructions operate on the raw bit representation of their operands.
This behaviour is the same for general-purpose and floating-point registers;
floating-point register contents are not interpreted as IEEE-754 values for
the purposes of logical operations or flag generation.

For a logical operation:

\begin{itemize}
\item \flagN\ is set from the most-significant bit of the result:
bit 31 for a general-purpose register and bit 63 for a floating-point register.
\item \flagZ\ is set if all bits of the result are zero.
\item \flagC\ and \flagV\ are cleared.
\item \flagE\ is cleared.
\item \flagI, \flagU, and \flagX\ are cleared.
\end{itemize}

\subsection{\code{and, andd}}

\begin{itemize}
\item Syntax: \code{and \$rd, \$ra, \$rb | andd \#rd, \#ra, \#rb}
\item Operation: \code{rd~\assign~ra~\&~rb}
\end{itemize}

\subsection{\code{andi}}

\begin{itemize}
\item Syntax: \code{andi \$rd, \$ra, imm16}
\item Operation: \code{rd~\assign~ra~\&~ZEXT(imm16)}
\end{itemize}

\subsection{\code{or, ord}}

\begin{itemize}
\item Syntax: \code{or \$rd, \$ra, \$rb | ord \#rd, \#ra, \#rb}
\item Operation: \code{rd~\assign~ra~|~rb}
\end{itemize}

\subsection{\code{ori}}

\begin{itemize}
\item Syntax: \code{ori \$rd, \$ra, imm16}
\item Operation: \code{rd~\assign~ra~|~ZEXT(imm16)}
\end{itemize}

\subsection{\code{xor, xord}}

\begin{itemize}
\item Syntax: \code{xor \$rd, \$ra, \$rb | xord \#rd, \#ra, \#rb}
\item Operation: \code{rd~\assign~ra~\^~rb}
\end{itemize}

\subsection{\code{xori}}

\begin{itemize}
\item Syntax: \code{xori \$rd, \$ra, imm16}
\item Operation: \code{rd~\assign~ra~\^~ZEXT(imm16)}
\end{itemize}

\subsection{\code{not, notd}}

\begin{itemize}
\item Syntax: \code{not \$rd, \$ra | notd \#rd, \#ra}
\item Operation: \code{rd~\assign~\~~ra}
\end{itemize}

\section{Comparison Instructions}

Floating-point comparisons use IEEE-754 binary64 numerical comparison semantics.

If either floating-point operand is \code{NaN}, the comparison is unordered.
For a flag-setting floating-point comparison, an unordered comparison sets
\flagE\ and clears the ordered comparison flags. No ordering relationship is
considered true for an unordered comparison.

For direct floating-point register-comparison branches, unordered comparison
behaviour is defined by the branch condition itself: \code{bd.eq}, \code{bd.lt},
\code{bd.le}, \code{bd.gt}, and \code{bd.ge} do not branch, while \code{bd.ne} branches.

Floating-point comparison does not compare IEEE-754 encodings bitwise; values are
compared numerically. In particular, \code{+0.0} and \code{-0.0} compare equal.

\subsection{\code{cmp, cmpd}}

\begin{itemize}
  \item Syntax: \code{cmp \$ra, \$rb | cmpd \#ra, \#rb}
  \item Operation: compares \code{ra} with \code{rb} and updates the condition
        flags without storing a result.
  \item \code{cmp} performs a 32-bit integer comparison. The comparison updates
        the integer condition flags as though \code{ra - rb} had been evaluated,
        but not arithmetic result is written to a register.
  \item \code{cmpd} performs an IEEE 754 binary64 numerical comparison.
        \code{Z} is set when the operands compare equal and \code{N} is set when
        \code{ra} compares less than \code{rb}. \code{C} and \code{V} are cleared.
  \item If either operand of \code{cmpd} is \code{NaN}, the comparison is unordered.
        \code{E} is set and the ordered comparison flags are cleared.
  \item Only \code{\%sr} is modified.
\end{itemize}

\subsection{\code{cmpi}}

\begin{itemize}
\item Syntax: \code{cmpi \$ra, imm16}
\item Operation: computes \code{ra~--~SEXT(imm16)} for flag results only.
\item Only \code{\%sr} is modified.
\end{itemize}

\subsection{btst, btstd}

\begin{itemize}
\item Syntax: \code{btst \$ra, imm14 | btstd \#ra, imm14}
\item Operation: test bit \code{imm14} of the raw bit representation of \code{ra}.
\item For \code{btst}, valid bit indices are \code{0--31}.
\item For \code{btstd}, valid bit indices are \code{0--63}.
\item On a successful bit test:
      \begin{itemize}
      \item Z is set if the selected bit is clear, and cleared if the selected bit is set.
      \item N, C, V, E, I, U, and X are cleared.
      \end{itemize}
      \item If \code{imm14} is outside the valid bit-index range:
      \begin{itemize}
      \item E is set.
      \item N, Z, C, V, I, U, and X are cleared.
      \end{itemize}
\item No general-purpose or floating-point register is modified.
\end{itemize}

\section{Shift Instructions}

Shift instructions that use a register to hold the shift value always use a
general purpose register. These instructions do not take the underlying format
into account, so while their use on floating-point registers is nonsensical from
a floating-point perspective, they open up the possibility of 64-bit integer
manipulations for those who need such functionality.

\subsection{Shift Register Behaviour}

Shift instructions operate on the raw bit representation of the source
register. This behaviour is the same for general-purpose and floating-point
registers; floating-point register contents are not interpreted as IEEE-754
values for the purposes of shift operations or flag generation.

For a successful shift operation:

\begin{itemize}
\item \flagN\ is set from the most-significant bit of the result:
bit 31 for a general-purpose register and bit 63 for a floating-point register.
\item \flagZ\ is set if all bits of the result are zero.
\item \flagC\ contains the last bit shifted out of the register.
\item \flagV\ and \flagE\ are cleared.
\item \flagI, \flagU, and \flagX\ are cleared.
\end{itemize}

If the shift count is greater than or equal to the bit-width of the target register:

\begin{itemize}
\item \flagE\ is set.
\item \flagN, \flagZ, \flagC, \flagV, \flagI, \flagU, and \flagX\ are cleared.
\item The destination register is left unchanged.
\end{itemize}

\subsection{\code{shl, shld}}

\begin{itemize}
\item Syntax: \code{shl \$rd, \$ra, \$rb | shld \#rd, \#ra, \$rb }
\item Operation: \code{rd~\assign~ra~<{}<~rb}
\item Vacated low bits are filled with zero.
\end{itemize}

\subsection{\code{shli, shldi}}

\begin{itemize}
\item Syntax: \code{shli \$rd, \$ra, imm14 | shldi \#rd, \#ra, imm14}
\item Operation: \code{rd~\assign~ra~<{}<~imm14}
\item Vacated low bits are filled with zero.
\end{itemize}

\subsection{\code{shr, shrd}}

\begin{itemize}
\item Syntax: \code{shr \$rd, \$ra, \$rb | shrd \#rd, \#ra, \$rb}
\item Operation: logical right shift, \code{rd~\assign~ra~>{}>~rb}
\item Vacated high bits are filled with zero.
\end{itemize}

\subsection{\code{shri, shrdi}}

\begin{itemize}
\item Syntax: \code{shri \$rd, \$ra, imm14 | shrdi \#rd, \#ra, imm14}
\item Operation: logical right shift, \code{rd~\assign~ra~>{}>~imm14}
\item Vacated high bits are filled with zero.
\end{itemize}

\subsection{\code{sar, sard}}

\begin{itemize}
\item Syntax: \code{sar \$rd, \$ra, \$rb | sard \#rd, \#ra, \$rb}
\item Operation: arithmetic right shift, \code{rd~\assign~ra~>{}>~rb}
\item Vacated high bits are filled with the original sign bit of (ra).
\end{itemize}

\subsection{\code{sari, sardi}}

\begin{itemize}
\item Syntax: \code{sari \$rd, \$ra, imm14 | sardi \#rd, \#ra, imm14}
\item Operation: arithmetic right shift, \code{rd~\assign~ra~>{}>~imm14}
\item Vacated high bits are filled with the original sign bit of (ra).
\end{itemize}

\section{Bit Manipulation Instructions}

For a successful bit-manipulation operation:

\begin{itemize}
\item N is set from the most-significant bit of the result: bit 31 for a
      general-purpose register and bit 63 for a floating-point register.
\item Z is set if the resulting register value is zero.
\item C, V, E, I, U, and X are cleared.
\end{itemize}

If the selected bit index is outside the valid range of the target register:

\begin{itemize}
\item E is set.
\item N, Z, C, V, I, U, and X are cleared.
\item The destination register is left unchanged.
\end{itemize}

\subsection{\code{bset, bsetd}}

\begin{itemize}
\item Syntax: \code{bset \$rd, \$ra, imm14 | bsetd \#rd, \#ra, imm14}
\item Operation: sets \code{ra.imm14}, then writes the result to \code{rd}.
\end{itemize}

\subsection{\code{bclr, bclrd}}

\begin{itemize}
\item Syntax: \code{bclr \$rd, \$ra, imm14 | bclrd \#rd, \#ra, imm14}
\item Operation: clears \code{ra.imm14}, then writes the result to \code{rd}.
\end{itemize}

\subsection{\code{btgl, btgld}}

\begin{itemize}
\item Syntax: \code{btgl \$rd, \$ra, imm14 | btgld \#rd, \#ra, imm14}
\item Operation: toggles bit \code{ra.imm14} of \code{ra}, then writes the
      result to \code{rd}.
\end{itemize}

\section{Multiply and Divide Instructions}

Multiply and divide instructions on general purpose registers utilise two
target registers, to hold upper and lower result and quotient and remainder
respectively. Multiply and divide instructions on floating-point registers
only require a single target register.

\subsection{\code{mul}}

\begin{itemize}
  \item Syntax: \code{mul \$rd0, \$rd1, \$ra, \$rb}
  \item Operation: signed 32-bit multiply, \code{rd1:rd0~\assign~ra~$\times$~rb}
  \item The operands \code{ra} and \code{rb} are interpreted as signed 32-bit
        two's complement integers.
  \item The full signed 64-bit product is computed.
  \item The low 32 bits of the product are written to \code{rd0}.
  \item The high 32 bits of the product are written to \code{rd1}.
  \item On successful multiplication:
  \begin{itemize}
  \item \flagN\ is set from bit 31 of \code{rd1}.
  \item \flagZ\ is set if the full 64-bit product is zero.
  \item Flags \flagC, \flagV\ and \flagE\ are cleared.
  \end{itemize}
\end{itemize}

\subsection{\code{muld}}

\begin{itemize}
\item Syntax: \code{muld \#rd, \#ra, \#rb}
\item Operation: IEEE-754 binary64 multiplication,
      \code{rd~\assign~ra~$\times$~rb}.
\item Updates flags:
      \begin{itemize}
      \item \flagN\ is set when the numeric result is negative.
            \code{+0.0} and \code{-0.0} do not set \flagN.
      \item \flagZ\ is set if the result is \code{+0.0} or \code{-0.0}.
      \item \flagC\ and \flagV\ are cleared.
      \item \flagI\ is set if the result is positive or negative \code{infinity}.
      \item \flagU\ is set if the result underflows binary64 representation.
      \item \flagX\ is set if the exact mathematical result cannot be represented
            exactly in binary64 and rounding is required.
      \item \flagE\ is set if the operation produces an invalid floating-point
          result, including \code{NaN}; otherwise it is cleared.
      \end{itemize}
\end{itemize}

\subsection{\code{mulu}}

\begin{itemize}
  \item Syntax: \code{mulu \$rd0, \$rd1, \$ra, \$rb}
  \item Operation: unsigned 32-bit multiply,
        \code{rd1:rd0~\assign~ra~$\times$~rb}
  \item The operands \code{ra} and \code{rb} are interpreted as unsigned 32-bit
        integers.
  \item The full unsigned 64-bit product is computed.
  \item The low 32 bits of the product are written to \code{rd0}.
  \item The high 32 bits of the product are written to \code{rd1}.
  \item On successful multiplication:
  \begin{itemize}
  \item \flagN\ is set from bit 31 of \code{rd1}.
  \item \flagZ\ is set if the full 64-bit product is zero.
  \item Flags \flagC\ \flagV\ and \flagE\ are cleared.
  \end{itemize}
\end{itemize}

\subsection{\code{div}}

\begin{itemize}
\item Syntax: \code{div \$rd0, \$rd1, \$ra, \$rb}
\item Operation: signed 32-bit division.
\item The operands \code{ra} and \code{rb} are interpreted as signed 32-bit
two's-complement integers.
\item On successful division, the signed quotient is written to \code{rd0}
and the signed remainder is written to \code{rd1}.

\item If \code{rb} is zero, division cannot be performed:
\begin{itemize}
\item \flagE\ is set.
\item \flagN, \flagZ, \flagC, and \flagV\ are cleared.
\item \code{rd0} and \code{rd1} are left unchanged.
\end{itemize}

\item If \code{ra} is \code{0x80000000} and \code{rb} is
\code{0xFFFFFFFF}, the mathematical quotient cannot be represented as
a signed 32-bit integer:
\begin{itemize}
\item \flagV\ and \flagE\ are set.
\item \flagN, \flagZ, and \flagC\ are cleared.
\item \code{rd0} and \code{rd1} are left unchanged.
\end{itemize}

\item Division errors do not trap in VIV-32 v\version{}.

\item On successful division:
\begin{itemize}
\item \flagN\ is set if the quotient is negative.
\item \flagZ\ is set if the quotient is zero.
\item \flagC\ is set if the remainder is non-zero.
\item \flagV\ and \flagE\ are cleared.
\end{itemize}

\item As an integer operation, \flagI, \flagU, and \flagX\ are cleared.
\end{itemize}

\subsection{\code{divd}}

\begin{itemize}
\item Syntax: \code{divd \#rd, \#ra, \#rb}
\item Operation: IEEE-754 binary64 division,
      \code{rd~\assign~ra~/~rb}.
\item Updates flags:
      \begin{itemize}
      \item \flagN\ is set when the numeric result is negative.
            \code{+0.0} and \code{-0.0} do not set \flagN.
      \item \flagZ\ is set if the result is \code{+0.0} or \code{-0.0}.
      \item \flagC\ and \flagV\ are cleared.
      \item \flagI\ is set if the result is positive or negative \code{infinity}.
      \item \flagU\ is set if the result underflows binary64 representation.
      \item \flagX\ is set if the exact mathematical result cannot be represented
            exactly in binary64 and rounding is required.
      \item \flagE\ is set if the operation encounters an invalid or exceptional
            arithmetic condition, including division by zero or production of
            \code{NaN}; otherwise it is cleared.
      \end{itemize}
\end{itemize}

\subsection{\code{divu}}

\begin{itemize}
\item Syntax: \code{divu \$rd0, \$rd1, \$ra, \$rb}
\item Operation: unsigned 32-bit division.
\item The operands \code{ra} and \code{rb} are interpreted as unsigned
      32-bit integers.
\item On successful division, the unsigned quotient is written to
      \code{rd0} and the unsigned remainder is written to \code{rd1}.
\item If \code{rb} is zero, division cannot be performed:
      \begin{itemize}
      \item \flagE\ is set.
      \item \flagN, \flagZ, \flagC, and \flagV\ are cleared.
      \item \code{rd0} and \code{rd1} are left unchanged.
      \end{itemize}
\item Divide-by-zero does not trap in VIV-32 v\version{}.
\item On successful division:
      \begin{itemize}
      \item \flagN\ is set from bit 31 of the quotient.
      \item \flagZ\ is set if the quotient is zero.
      \item \flagC\ is set if the remainder is non-zero.
      \item \flagV\ and \flagE\ are cleared.
      \end{itemize}
\item As an integer operation, \flagI, \flagU, and \flagX\ are cleared.
\end{itemize}

\section{Load and Store Instructions}

Load and store instructions transfer data between registers and
memory. VIV-32 uses base-plus-offset addressing for ordinary memory access.
Multi-byte load and store operations are big-endian.

\subsection{Address Calculation}
\label{sec:load-store-address-calculation}

The assembly syntax for memory addressing is:

\begin{center}
\code{reg, [\$base, offset15]}
\end{center}

\noindent
In this syntax, \code{base} is a general-purpose register containing a 32-bit
base address, and \code{offset15} is sign-extended to 32-bits. \code{reg} may
be either a general purpose or floating-point register, depending on the
encoded instruction.

\noindent
The effective address is computed as:

\begin{center}
\code{address~\assign~[base~+~SEXT(offset15)]}
\end{center}

\noindent
The computed address is interpreted according to the access size of the
instruction. Byte accesses may occur at any address. Halfword, word and
double-half accesses are subject to the alignment rules defined below.

\subsection{Alignment}

\begin{itemize}
\item Byte loads and stores may access any address.
\item Halfword loads and stores must access a 2-byte aligned address.
\item Word and double-half loads and stores must access a 4-byte aligned
      address.
\item Misaligned halfword, word and double-half accesses raise a
      \code{MisalignedAccessException}:
      \begin{itemize}
      \item \code{\%ecause} receives the \code{MisalignedAccessException}.
      \item The exception entry sequence is defined in the exceptions and
            interrupts chapter.
      \end{itemize}
\end{itemize}

\subsection{Flag Behaviour}

\begin{itemize}
\item Load instructions update condition flags according to the resulting
destination-register value.

\item Integer loads update the integer result flags:
\begin{itemize}
\item \flagN\ reflects the sign of the loaded 32-bit result.
\item \flagZ\ is set when the loaded result is zero.
\item \flagC, \flagV, \flagE, \flagI, \flagU, and \flagX\ are cleared.
\end{itemize}

\item Floating-point loads update the floating-point result flags:
\begin{itemize}
\item \flagN\ is set when the resulting floating-point value is negative.
\item \flagZ\ is set when the result is \code{+0.0} or \code{-0.0}.
\item \flagI\ is set when the resulting value is positive or negative
infinity.
\item \flagU, \flagX, and \flagE\ are updated according to any conversion
performed by the load.
\item \flagC\ and \flagV\ are cleared.
\end{itemize}

\item The binary64 half-load instructions \code{ldl} and \code{ldu} update the
flags from the complete 64-bit contents of the destination floating-point
register after the selected half has been written.

\item When a binary64 value is assembled using two half-loads, the flag state
produced by the first half-load is generally not meaningful. The flags
describe the complete loaded value only after both halves have been
loaded.

\item Store instructions do not update condition flags.
\end{itemize}

\subsection{\code{lb}}

\begin{itemize}
\item Syntax: \code{lb \$rd, [\$base, offset]}
\item Operation: load signed byte.
\item One byte is read from memory, sign-extended to 32 bits, then put in
      general purpose register \code{\$rd}.
\end{itemize}

\subsection{\code{lbu}}

\begin{itemize}
\item Syntax: \code{lbu \$rd, [\$base, offset]}
\item Operation: load unsigned byte.
\item One byte is read from memory, zero-extended to 32 bits, then put in
      general purpose register \code{\$rd}.
\end{itemize}

\subsection{\code{lh}}

\begin{itemize}
\item Syntax: \code{lh \$rd, [\$base, offset]}
\item Operation: load signed halfword.
\item Two bytes are read from memory, sign-extended to 32-bits, then put in
      general purpose register \code{\$rd}.
\end{itemize}

\subsection{\code{lhu}}

\begin{itemize}
\item Syntax: \code{lhu \$rd, [\$base, offset]}
\item Operation: load unsigned halfword.
\item Two bytes are read from memory, zero-extended to 32-bits, then put in
      general purpose register \code{\$rd}.
\end{itemize}

\subsection{\code{lw}}

\begin{itemize}
\item Syntax: \code{lw \$rd, [\$base, offset]}
\item Operation: load word.
\item Four bytes are read from memory, then put in general purpose register
      \code{\$rd}.
\end{itemize}

\subsection{\code{ld}}

\begin{itemize}
\item Syntax: \code{ld \#rd, [\$base, offset]}
\item Operation: loads float.
\item Four bytes are read from memory, interpeted as an IEEE 754 binary32 type,
      extended to binary64, then placed in floating-point register  \code{\#rd}.
\end{itemize}

\subsection{\code{ldl}}

\begin{itemize}
\item Syntax: \code{ldl \#rd, [\$base, offset]}
\item Operation: load lower double-half.
\item Four bytes are read from memory, then put in the lower 32-bits of
      floating-point register \code{\#rd}. This operation is raw, leaving the
      upper bits untouched.
\item Condition flags are updated for \code{\#rd} as a whole.
\end{itemize}

\subsection{\code{ldu}}

\begin{itemize}
\item Syntax: \code{ldu \#rd, [\$base, offset]}
\item Operation: load upper double-half.
\item Four bytes are read from memory, then put in the upper 32-bits of
      floating-point register \code{\#rd}. This operation is raw, leaving the
      lower bits untouched.
\item Condition flags are updated for \code{\#rd} as a whole.
\end{itemize}

\subsection{\code{sb}}

\begin{itemize}
\item Syntax: \code{sb \$rs, [\$base, offset]}
\item Operation: store byte.
\item The low 8 bits of \code{\$rs} are written to the calculated memory
      address.
\end{itemize}

\subsection{\code{sh}}

\begin{itemize}
\item Syntax: \code{sh \$rs, [\$base, offset]}
\item Operation: store halfword.
\item The low 16 bits of \code{\$rs} are written to the calculated memory
      address memory using big-endian byte order.
\end{itemize}

\subsection{\code{sw}}

\begin{itemize}
\item Syntax: \code{sw \$rs, [\$base, offset]}
\item Operation: store word.
\item The full 32-bit value in \code{\$rs} is written to the calculated memory
      address using big-endian byte order.
\end{itemize}

\subsection{\code{sd}}

\begin{itemize}
\item Syntax: \code{sd \#rs, [\$base, offset]}
\item Operation: store an IEEE-754 binary32 floating-point value.
\item The binary64 value held in floating-point register \code{\#rs} is
      converted to IEEE-754 binary32 and the resulting 32-bit representation
      is written to the calculated memory address using big-endian byte order.
\item The conversion performed by \code{sd} is blind: \code{sd} does not
      report whether precision or range was lost during conversion and does not
      update condition flags.
\item Software that requires assurance that the binary64 source can be
      represented appropriately as binary32 should execute \code{disf} before
      performing the store.
\end{itemize}

\subsection{\code{sdl}}

\begin{itemize}
\item Syntax: \code{sdl \#rs, [\$base, offset]}
\item Operation: store lower binary64.
\item The lower 32-bit value from floating-point register \code{\#rs} is written
      to the calculated memory address using big-endian byte order.
\end{itemize}

\subsection{\code{sdu}}

\begin{itemize}
\item Syntax: \code{sdu \#rs, [\$base, offset]}
\item Operation: store upper binary64.
\item The upper 32-bit value from floating-point register \code{\#rs} is written
      to the calculated memory address using big-endian byte order.
\end{itemize}

\section{Constant Construction Instructions}

VIV-32 instructions are fixed-width 32-bit words. A single instruction cannot
encode every possible 32-bit constant together with an opcode and register
fields. Larger constants and addresses are therefore constructed using one or
more instructions.

Constants construction instructions set the conditions register \code{\%sr} as
per normal.

\subsection{\code{lui}}

\begin{itemize}
  \item Syntax: \code{lui \$rd, imm16}
  \item Operation: load upper immediate, \code{rd~\assign~imm16~<{}<~16}
  \item The 16-bit immediate is placed in bits 31--16 of \code{\$rd}, with bits
        15--0 of the result set to zero.
\end{itemize}

\subsection{\code{lli}}

\begin{itemize}
  \item Syntax: \code{lli \$rd, imm16}
  \item Operation: load lower immediate, \code{rd~\assign~imm16~\&~0xFFFF}
  \item The 16-bit immediate is zero-extended and written to \code{\$rd}.
\end{itemize}

\subsection{\code{lhi}}

\begin{itemize}
  \item Syntax: \code{lhi \$rd, imm16}
  \item Operation: load upper immediate, preserving lower bytes,

        \code{\$rd~\assign~(\$rd~\&~0xFFFF)~|~(imm16~<{}<~16)}
\end{itemize}

\subsection{Full-Width Constants, \code{li} and \code{la}}

The assembler provides pseudo instructions \code{li} and \code{la} to load
32-bit values and addresses respectively.

\section{Inter-File Transfer Instructions}

\subsection{Floating-Point Rounding}

VIV-32 v0.2 does not provide a programmable floating-point rounding-mode
control register.

Unless explicitly specified otherwise by an instruction, floating-point
arithmetic and floating-point conversions use IEEE-754 round-to-nearest,
ties-to-even.

This rounding mode applies to operations including:

\begin{itemize}
\item IEEE-754 binary64 arithmetic;
\item binary64-to-binary32 conversion;
\item integer-to-binary64 conversion where rounding is required;
\item binary64-to-integer conversion.
\end{itemize}

Where a conversion or arithmetic operation cannot represent the exact
mathematical result in its destination format, \flagX\ is set when rounding
is required, subject to the instruction-specific flag rules.

\subsection{\code{disf}}

\begin{itemize}
\item Syntax: \code{disf \#ra}
\item Operation: tests the result of converting the IEEE-754 binary64 value
      in \code{\#ra} to IEEE-754 binary32 without storing the converted value.
\item The source floating-point register is not modified.
\item The condition flags describe the binary32 conversion that would be
      performed by \code{dtof} or \code{sd}:
      \begin{itemize}
      \item \flagN\ is set if the converted value is numerically negative.
      \item \flagZ\ is set if the converted value is \code{+0.0} or \code{-0.0}.
      \item \flagI\ is set if the converted value is positive or negative
            \code{infinity}.
      \item \flagU\ is set if conversion enters the binary32 underflow range.
      \item \flagX\ is set if conversion to binary32 requires rounding or
            otherwise loses precision.
      \item \flagE\ is set only if the conversion is invalid.
      \item \flagC\ and \flagV\ are cleared.
      \end{itemize}
\item IEEE-754 \code{NaN} values are considered representable binary32
      values and may therefore be tested by \code{disf}; the presence of a
      \code{NaN} does not by itself make the conversion invalid.
\item No converted value is written to a register.
\end{itemize}

\subsection{\code{utod}}

\begin{itemize}
  \item Syntax: \code{utod \#rd, \$rs}
  \item Operation: \code{\#rd~\assign~FLOAT(\$rs)}
  \item Loads numerical value of unsigned general-purpose register
        \code{\$rs} into floating-point register \code{\#rd} as binary64,
        maintaining numerical value.
\end{itemize}

\subsection{\code{itod}}

\begin{itemize}
  \item Syntax: \code{itod \#rd, \$rs}
  \item Operation: \code{\#rd~\assign~FLOAT(\$rs)}
  \item Loads numerical value of signed general-purpose register
        \code{\$rs} into floating-point register \code{\#rd} as binary64,
        maintaining numerical value.
\end{itemize}

\subsection{\code{ftod}}

\begin{itemize}
  \item Syntax: \code{ftod \#rd, \$rs}
  \item Operation: \code{\#rd~\assign~BINARY64(\$rs)}
  \item Loads binary32 value of general-purpose register
        \code{\$rs} into floating-point register \code{\#rd} as binary64.
\end{itemize}

\subsection{\code{dtou}}

\begin{itemize}
\item Syntax: \code{dtou \$rd, \#rs}
\item Operation: \code{\$rd~\assign~UNSIGNED(\#rs)}
\item Converts the IEEE-754 binary64 numerical value in \code{\#rs} to an
      unsigned 32-bit integer using the architectural rounding mode.
\item On successful conversion:
      \begin{itemize}
      \item The converted value is written to \code{\$rd}.
      \item \flagN\ is set from bit 31 of the resulting integer.
      \item \flagZ\ is set if the resulting integer is zero.
      \item \flagC, \flagV, and \flagE\ are cleared.
      \item \flagI\ and \flagU\ are cleared.
      \item \flagX\ is set if rounding was required to produce the integer result;
      otherwise it is cleared.
      \end{itemize}
\item If the source is \code{NaN}, positive or negative infinity, or cannot
      be represented as an unsigned 32-bit integer after rounding:
      \begin{itemize}
      \item \flagE\ is set.
      \item \flagN, \flagZ, \flagC, \flagV, \flagI, \flagU, and \flagX\ are cleared.
      \item \code{\$rd} is left unchanged.
      \end{itemize}
\end{itemize}

\subsection{\code{dtoi}}

\begin{itemize}
\item Syntax: \code{dtoi \$rd, \#rs}
\item Operation: \code{\$rd~\assign~SIGNED(\#rs)}
\item Converts the IEEE-754 binary64 numerical value in \code{\#rs} to a
      signed 32-bit two's-complement integer using the architectural
      rounding mode.
\item On successful conversion:
      \begin{itemize}
      \item The converted value is written to \code{\$rd}.
      \item \flagN\ is set if the resulting integer is negative.
      \item \flagZ\ is set if the resulting integer is zero.
      \item \flagC, \flagV, and \flagE\ are cleared.
      \item \flagI\ and \flagU\ are cleared.
      \item \flagX\ is set if rounding was required to produce the integer result;
      otherwise it is cleared.
      \end{itemize}
\item If the source is \code{NaN}, positive or negative infinity, or cannot
be represented as a signed 32-bit integer after rounding:
      \begin{itemize}
      \item \flagE\ is set.
      \item \flagN, \flagZ, \flagC, \flagV, \flagI, \flagU, and \flagX\ are cleared.
      \item \code{\$rd} is left unchanged.
      \end{itemize}
\end{itemize}

\subsection{\code{dtof}}

\begin{itemize}
\item Syntax: \code{dtof \$rd, \#rs}
\item Operation: \code{\$rd~\assign~BINARY32(\#rs)}
\item Converts the IEEE-754 binary64 value in \code{\#rs} to IEEE-754
      binary32 using the architectural rounding mode and writes the resulting
      32-bit representation to \code{\$rd}.
\item The condition flags describe the resulting binary32 value:
      \begin{itemize}
      \item \flagN\ is set if the converted value is numerically negative.
      \item \flagZ\ is set if the converted value is \code{+0.0} or \code{-0.0}.
      \item \flagI\ is set if the converted value is positive or negative infinity.
      \item \flagU\ is set if conversion enters the binary32 underflow range.
      \item \flagX\ is set if conversion requires rounding or otherwise loses precision.
      \item \flagE\ is set only if the conversion is invalid.
      \item \flagC\ and \flagV\ are cleared.
      \end{itemize}
\item IEEE-754 \code{NaN} values are representable as binary32 values and do
      not by themselves make the conversion invalid.
\item If the conversion is invalid, \code{\$rd} is left unchanged.
\end{itemize}

\section{Branch Instructions}

Branch instructions perform conditional PC-relative control transfer. VIV-32
defines two branch families:

\begin{itemize}
  \item register-comparison branches, which compare two registers directly;
  \item flag-based branches, which test the existing condition flags in
  \code{\%sr}.
\end{itemize}

\subsection{Branch Target Calculation}

\begin{itemize}
  \item Branch offsets are signed PC-relative offsets.
  \item Branch offsets are measured in 32-bit instruction words.
  \item The encoded offset is shifted left by two bits to form a byte offset.
  \item The branch target is relative to the address of the instruction
        following the branch.
  \item Branch targets must be 4-byte aligned.
  \item Branch instructions do not update condition flags.
\end{itemize}

\subsection{Register-Comparison Branches}

Register-comparison branches compare \code{ra} and \code{rb} directly. They do
not depend on the current condition flags, and they do not update the condition
flags. The jump target is calculated in instruction words:

\code{target = \%pc + 4 + SEXT(imm14) <{}< 2}

\noindent
That is, the target is relative to the NEXT instruction, not the executed
instruction.
\begin{itemize}
\item \code{imm14} can be a label or a value.
\item Unsigned comparisons are only valid for general purpose registers.
\item Floating-point comparisons use IEEE-754 binary64 numerical comparison semantics.
\item If either floating-point operand is \code{NaN}, the comparison is unordered.
For unordered comparisons, \code{bd.eq}, \code{bd.lt}, \code{bd.le},
\code{bd.gt}, and \code{bd.ge} do not branch; \code{bd.ne} branches.
\item Register-comparison branches do not modify the condition flags, including for
unordered floating-point comparisons.
\end{itemize}

\subsection{\code{b.eq, bd.eq}}

\begin{itemize}
  \item Syntax: \code{b.eq \$ra, \$rb, imm14 | bd.eq \#ra, \#rb, imm14}
  \item Operation: branch if \code{ra~=~rb}.
  \item For general purpose registers, this is bitwise equivalence. For
        floating-point register, this is numerical equivalence.
\end{itemize}

\subsection{\code{b.ne, bd.ne}}

\begin{itemize}
  \item Syntax: \code{b.ne \$ra, \$rb, imm14 | bd.ne \#ra, \#rb, imm14}
  \item Operation: branch if \code{ra~$\ne$~rb}.
  \item For general purpose registers, this is bitwise inequivalence. For
        floating-point register, this is numerical inequivalence.
\end{itemize}

\subsection{\code{b.lt, bd.lt}}

\begin{itemize}
  \item Syntax: \code{b.lt \$ra, \$rb, imm14 | bd.lt \#ra, \#rb, imm14}
  \item Operation: branch if \code{ra~<~rb}, treating both operands as signed
        numbers.
\end{itemize}

\subsection{\code{b.le, bd.le}}

\begin{itemize}
  \item Syntax: \code{b.le \$ra, \$rb, imm14 | bd.le \#ra, \#rb, imm14}
  \item Operation: branch if \code{ra~$\le$~rb}, treating both operands as signed
        numbers.
\end{itemize}

\subsection{\code{b.gt, bd.gt}}

\begin{itemize}
  \item Syntax: \code{b.gt \$ra, \$rb, imm14 | bd.gt \#ra, \#rb, imm14}
  \item Operation: branch if \code{ra~>~rb}, treating both operands as signed
        numbers.
\end{itemize}

\subsection{\code{b.ge, bd.ge}}

\begin{itemize}
  \item Syntax: \code{b.ge \$ra, \$rb, imm14 | bd.ge \#ra, \#rb, imm14}
  \item Operation: branch if \code{ra~$\ge$~rb}, treating both operands as signed
        numbers.
\end{itemize}

\subsection{\code{b.ltu}}

\begin{itemize}
  \item Syntax: \code{b.ltu \$ra, \$rb, imm14}
  \item Operation: branch if \code{ra~<~rb}, treating both operands as unsigned
        32-bit integers.
\end{itemize}

\subsection{\code{b.leu}}

\begin{itemize}
  \item Syntax: \code{b.leu \$ra, \$rb, imm14}
  \item Operation: branch if \code{ra~$\le$~rb}, treating both operands as
        unsigned 32-bit integers.
\end{itemize}

\subsection{\code{b.gtu}}

\begin{itemize}
  \item Syntax: \code{b.gtu \$ra, \$rb, imm14}
  \item Operation: branch if \code{ra~>~rb}, treating both operands as
        unsigned 32-bit integers.
\end{itemize}

\subsection{\code{b.geu}}

\begin{itemize}
  \item Syntax: \code{b.geu \$ra, \$rb, imm14}
  \item Operation: branch if \code{ra~$\ge$~rb}, treating both operands as
        unsigned 32-bit integers.
\end{itemize}

\subsection{Flag-Based Branches}

Flag-based branches test a single current condition flag in the status register
\code{\%sr}. They do not perform a comparison, do not interpret the flags as
signed, unsigned, or floating-point state, and do not update any flags.

The jump target is calculated in instruction words:

\code{target = \%pc + 4 + SEXT(imm22) <{}< 2}

\noindent
That is, the target is relative to the instruction following the branch.
\code{imm22} may be a label or a value.

\subsection{\code{bf.ns}}

\begin{itemize}
\item Syntax: \code{bf.ns imm22}
\item Operation: branch if \flagN\ is set.
\end{itemize}

\subsection{\code{bf.nc}}

\begin{itemize}
\item Syntax: \code{bf.nc imm22}
\item Operation: branch if \flagN\ is clear.
\end{itemize}

\subsection{\code{bf.zs}}

\begin{itemize}
\item Syntax: \code{bf.zs imm22}
\item Operation: branch if \flagZ\ is set.
\end{itemize}

\subsection{\code{bf.zc}}

\begin{itemize}
\item Syntax: \code{bf.zc imm22}
\item Operation: branch if \flagZ\ is clear.
\end{itemize}

\subsection{\code{bf.cs}}

\begin{itemize}
\item Syntax: \code{bf.cs imm22}
\item Operation: branch if \flagC\ is set.
\end{itemize}

\subsection{\code{bf.cc}}

\begin{itemize}
\item Syntax: \code{bf.cc imm22}
\item Operation: branch if \flagC\ is clear.
\end{itemize}

\subsection{\code{bf.vs}}

\begin{itemize}
\item Syntax: \code{bf.vs imm22}
\item Operation: branch if \flagV\ is set.
\end{itemize}

\subsection{\code{bf.vc}}

\begin{itemize}
\item Syntax: \code{bf.vc imm22}
\item Operation: branch if \flagV\ is clear.
\end{itemize}

\subsection{\code{bf.es}}

\begin{itemize}
\item Syntax: \code{bf.es imm22}
\item Operation: branch if \flagE\ is set.
\end{itemize}

\subsection{\code{bf.ec}}

\begin{itemize}
\item Syntax: \code{bf.ec imm22}
\item Operation: branch if \flagE\ is clear.
\end{itemize}

\subsection{\code{bf.is}}

\begin{itemize}
\item Syntax: \code{bf.is imm22}
\item Operation: branch if \flagI\ is set.
\end{itemize}

\subsection{\code{bf.ic}}

\begin{itemize}
\item Syntax: \code{bf.ic imm22}
\item Operation: branch if \flagI\ is clear.
\end{itemize}

\subsection{\code{bf.us}}

\begin{itemize}
\item Syntax: \code{bf.us imm22}
\item Operation: branch if \flagU\ is set.
\end{itemize}

\subsection{\code{bf.uc}}

\begin{itemize}
\item Syntax: \code{bf.uc imm22}
\item Operation: branch if \flagU\ is clear.
\end{itemize}

\subsection{\code{bf.xs}}

\begin{itemize}
\item Syntax: \code{bf.xs imm22}
\item Operation: branch if \flagX\ is set.
\end{itemize}

\subsection{\code{bf.xc}}

\begin{itemize}
\item Syntax: \code{bf.xc imm22}
\item Operation: branch if \flagX\ is clear.
\end{itemize}

\section{Jump and Call Instructions}

Jump and call instructions perform unconditional control transfer. Unlike branch
instructions, jump and call instructions do not test condition flags or compare
registers.

\subsection{Jump Target Rules}

\begin{itemize}
\item Jump and call targets are 32-bit instruction addresses.
\item The encoded offset is measured in 32-bit instruction words.
\item The encoded offset is shifted left by two bits to form a byte offset.
\item The target is relative to the address of the following instruction.
\item Jump and call targets must be 4-byte aligned. This is ensured by the
      calculation:

      \code{\%pc~\assign~\%pc + 4 + SEXT(imm26) <{}< 2}
\item Jump and call instructions do not update condition flags.
\item \code{imm26} can be a value or label.
\end{itemize}

\subsection{\code{jmp}}

\begin{itemize}
\item Syntax: \code{jmp imm26}
\item Operation: unconditional jump.
\end{itemize}

\subsection{\code{call}}

\begin{itemize}
\item Syntax: \code{call imm26}
\item Operation: unconditional jump with updated link register \code{\%lr}
      (\code{\$15}). The return address is the address of the instruction
      following the \code{call}.

      \code{\$15~\assign~\%pc + 4}
\end{itemize}

\subsection{\code{jr}}

\begin{itemize}
\item Syntax: \code{jr \$target}
\item Operation: register-indirect unconditional jump,
      \code{\%pc~\assign~\$target}
\item The \code{target} register contains the destination address.
\item The destination address must be 4-byte aligned. Causes a
      \code{MisalignedFetchException} if the \code{target} address is not
      word-aligned.
\end{itemize}

\subsection{\code{jalr}}

\begin{itemize}
\item Syntax: \code{jalr \$rd, \$target}
\item Operation: register-indirect call, \code{\%pc~\assign~UNSIGNED(target)}
\item The \code{target} register contains the destination address.
\item The destination address must be 4-byte aligned. Causes a
      \code{MisalignedFetchException} if the \code{target} address is not
      word-aligned.
\item The return address is written to \code{rd}, \code{rd~\assign~\%pc + 4}
\end{itemize}

\subsection{\code{ret}}

\begin{itemize}
\item \code{ret} is an assembler pseudo-instruction, not a distinct
architectural instruction in VIV-32 v\version.
\item Syntax: \code{ret}
\item Expansion: \code{jr \$15}
\item Operation: return to the address held in the link register.
\item No condition flags are updated.
\end{itemize}

\subsection{Nested Calls}

\begin{itemize}
\item A call instruction overwrites the link register.
\item A function that calls another function must save its incoming link
      register value if it needs to return to its own caller, e.g.:

      \code{push \$lr}

      \code{call func}

      \code{pop \$lr}
\end{itemize}

\section{System and Control Instructions}

System and control instructions provide access to processor control state,
exception handling, interrupt control, and execution control.

\subsection{Control Registers}

\begin{itemize}
  \item VIV-32 defines a small set of architectural control registers.
  \item Control registers are accessed using \code{mrs} and \code{msr}.
  \item Access permissions are defined by the privilege model. In VIV-32
        v\version, only machine mode is defined.
\end{itemize}

\begin{center}
\begin{tabular}{lll}
\textbf{Name} & \textbf{Meaning} & \textbf{Description} \\
\code{\%pc}     & Program counter           & Address of the currently executing instruction \\
\code{\%sr}     & Status register           & Condition flags and interrupt-enable state \\
\code{\%epc}    & Exception program counter & Continuation address saved on exception entry \\
\code{\%esr}    & Exception status register & Status register saved on exception entry \\
\code{\%ecause} & Exception cause           & Cause code saved on exception entry \\
\code{\%edata}  & Exception data            & Exception-specific data saved on exception entry \\
\code{\%evbase} & Exception vector base     & Base address for exception and interrupt vectors \\
\end{tabular}
\end{center}

\noindent
For the purposes of instruction semantics, \code{\%pc} contains the address of the instruction
currently being executed. References to \code{\%pc + 4} therefore refer to the address of the
following sequential instruction.

At an architectural instruction boundary, before the next instruction begins execution,
\code{\%pc} identifies the next instruction to be executed.

\subsection{\code{nop}}

\begin{itemize}
  \item Syntax: \code{nop}
  \item Operation: no operation.
  \item Encoding: \code{0x00000000}.
  \item No general-purpose register is modified.
  \item No control register is modified.
  \item No condition flags are updated.
  \item Execution continues at the next instruction.
\end{itemize}

\subsection{\code{halt}}

\begin{itemize}
  \item Syntax: \code{halt}
  \item Operation: halt instruction execution.
  \item The processor stops executing instructions.
  \item The processor remains halted until reset or platform-defined external action.
  \item No condition flags are updated.
\end{itemize}

\subsection{\code{trap}}

\begin{itemize}
  \item Syntax: \code{trap imm12}
  \item Operation: raise a software trap exception.
  \item The immediate operand is encoded as the X-type \code{imm12} payload.
  \item On trap entry, \code{\%ecause} receives \code{SoftwareTrap}.
  \item On trap entry, \code{\%edata} receives \code{ZEXT(imm12)}.
  \item On trap entry, \code{\%epc} receives the address of the instruction following
        the \code{trap} instruction.
  \item The exception entry sequence is defined in the exceptions and interrupts chapter.
  \item No condition flags are updated.
\end{itemize}

\subsection{\code{syscall}}

\begin{itemize}
  \item Syntax: \code{syscall}
  \item Operation: raise a system call exception.
  \item \code{syscall} is a distinct architectural instruction.
  \item \code{syscall} may share exception-entry machinery with \code{trap}.
  \item On syscall entry, \code{\%epc} receives the address of the instruction following the \code{syscall}.
  \item On syscall entry, \code{\%ecause} receives the system call exception code.
  \item The system call number and arguments are passed according to the ABI.
  \item The exception entry sequence is defined in the exceptions and interrupts chapter.
  \item No condition flags are updated.
\end{itemize}

\subsection{iret}

\begin{itemize}
\item Syntax: \code{iret}
\item Operation: return from an exception or interrupt.
\item The program counter is restored from \code{\%epc}.
\item The processor status register is restored from \code{\%esr}.
\item Operation:
\end{itemize}

$$
\code{\%pc} \leftarrow \code{\%epc}
$$

$$
\code{\%sr} \leftarrow \code{\%esr}
$$

\begin{itemize}
\item All condition flags and the interrupt-enable state are restored as part of the
\code{\%sr} restoration.
\end{itemize}

\subsection{\code{ei}}

\begin{itemize}
  \item Syntax: \code{ei}
  \item Operation: enable maskable interrupts.
  \item The interrupt-enable bit in \code{\%sr} is set.
  \item No general-purpose register is modified.
  \item No condition flags are updated.
\end{itemize}

\subsection{\code{di}}

\begin{itemize}
  \item Syntax: \code{di}
  \item Operation: disable maskable interrupts.
  \item The interrupt-enable bit in \code{\%sr} is cleared.
  \item No general-purpose register is modified.
  \item No condition flags are updated.
\end{itemize}

\subsection{\code{rdpc}}

\begin{itemize}
  \item Syntax: \code{rdpc \$rd}
  \item Operation: read program counter.
  \item The address of the \code{rdpc} instruction is written to \code{rd}.
  \item Operation: \code{rd~\assign~\%pc}
  \item No condition flags are updated.
\end{itemize}

\subsection{\code{mrs}}

\begin{itemize}
  \item Syntax: \code{mrs \$rd, \%cr}
  \item Operation: move from control register.
  \item The value of control register \code{cr} is written to \code{rd}.
  \item No condition flags are updated.
\end{itemize}

\subsection{\code{msr}}

\begin{itemize}
  \item Syntax: \code{msr \%cr, \$rs}
  \item Operation: move to control register.
  \item The value of \code{rs} is written to control register \code{cr}. Reserved
        bits are not updated.
  \item Writes to control registers may have immediate architectural effects.
  \item \code{\%pc} is read-only through the control-register interface.
  \item Executing \code{msr \%pc, \$rs} raises an
        \code{IllegalInstructionException}. The program counter is not modified by the
        control-register write.
  \item No condition flags are updated, except when writing \code{\%sr}.
\end{itemize}

\subsubsection{Writing to \code{\%evbase}:}
\begin{itemize}
\item Writing \code{\%evbase} requires the source value to be 16-byte aligned.
\item If the source value is not 16-byte aligned, \code{\%evbase} is left unchanged and
the processor enters the halted state.
\item This condition does not raise an exception.
\end{itemize}


\section{Pseudo-Instructions}

Pseudo-instructions are assembler conveniences. They do not necessarily correspond
to distinct architectural instruction encodings. A conforming assembler may expand
pseudo-instructions into one or more real VIV-32 instructions.

\subsection{Pseudo-Instruction Policy}

\begin{itemize}
  \item Pseudo-instructions are not part of the architectural instruction encoding
        unless explicitly stated otherwise.
  \item Pseudo-instructions may expand differently depending on immediate size,
        relocation requirements, or target address range.
  \item The assembler is responsible for selecting a correct expansion.
  \item The exact expansion rules are defined by the assembler and toolchain specification.
  \item Pseudo-instructions do not define independent architectural side effects beyond
        those produced by their expansion into real VIV-32 instructions.
  \item Unless a pseudo-instruction definition explicitly constrains the expansion,
        condition flags, control-register effects, exceptions, and other observable
        behaviour are those of the selected instruction sequence.
  \item A conforming assembler must preserve the functional meaning of the
        pseudo-instruction, but different valid expansions may differ in incidental
        architectural side effects where this specification does not require otherwise.
\end{itemize}

\subsection{\code{ldd}}

\begin{itemize}
  \item Syntax: \code{ldd \#rd, [\$base, imm15]}
  \item Operation: loads a binary64 value from the calculated memory address, as
        per the load/store memory address algorithm, into floating-point register
        \code{rd}. Loads in big-endian order, with the upper double-half at the
        calculated address, and the lower double-half \code{address + 4}.
\end{itemize}

\subsection{\code{sdd}}

\begin{itemize}
  \item Syntax: \code{sdd \#rs, [\$base, imm15]}
  \item Operation: stores a binary64 value from floating-point register
        \code{rs} into the calculated memory address, as per the load/store
        memory address algorithm. Stores in big-endian order, with the upper
        double-half at the calculated address, and the lower double-half at
        \code{address + 4}.
\end{itemize}


\subsection{\code{li}}

\begin{itemize}
  \item Syntax: \code{li \$rd, imm32}
  \item Operation: load a 32-bit immediate value into \code{rd}.
  \item The assembler expands \code{li} into one or more constant-construction
         instructions, e.g.:

        \code{lli rd, imm32 \& 0xFFFF}
        
        \code{lhi rd, (imm32 >{}> 16) \& 0xFFFF}
        
\end{itemize}

\subsection{\code{la}}

\begin{itemize}
  \item Syntax: \code{la \$rd, symbol}
  \item Operation: load the address of \code{symbol} into \code{rd}.
  \item \code{symbol} \textbf{MUST} be a label.
  \item The assembler expands \code{la} into relocation-aware address
        construction instructions. The final address will be resolved by the
        linker.
\end{itemize}

\subsection{\code{mov, movd}}

\begin{itemize}
  \item Syntax: \code{mov \$rd, \$rs | movd \#rd, \#rs}
  \item Operation: copy the value of \code{rs} into \code{rd}, \code{rd~\assign~rs}
  \item The internal construction ensures flags are updated accordingly.
\end{itemize}

\subsection{\code{clr, clrd}}

\begin{itemize}
  \item Syntax: \code{clr \$rd | clrd \#rd}
  \item Operation: write zero to \code{rd}, \code{rd~\assign~0}
  \item The internal construction ensure flags are updated accordingly.
\end{itemize}

\subsection{\code{ret}}

See definition under Jump and Call Instructions.

\subsection{\code{push, pushd}}

\begin{itemize}
\item Syntax: \code{push \$rs | pushd \#rs}
\item Operation: push a register value onto the stack.
\item The stack grows downward.
\item The assembler expands the pseudo-instruction into the appropriate
      stack-pointer adjustment and store instruction sequence.
\item Preservation of ABI stack-alignment requirements, including alignment
      at function-call boundaries, is the responsibility of the toolchain
      and generated instruction sequence rather than the architectural
      \code{push} or \code{pushd} pseudo-instruction itself.
\end{itemize}

\subsection{\code{pop, popd}}

\begin{itemize}
  \item Syntax: \code{pop \$rd | popd \#rd}
  \item Operation: pops a register value from the stack.
  \item Care must be taken by the programmer to ensure this is used in
        conjunction with the appropriate \code{push, pushd}. If the stack is
        not word-aligned, this will cause a \code{MisalignedAccessException}.
\end{itemize}

\section{Reserved Instructions}

\begin{itemize}
  \item Instruction encodings not defined by VIV-32 v\version{} are reserved.
  \item Reserved encodings must not be generated by conforming assemblers.
  \item Reserved encodings are reserved for future architectural extensions, implementation-specific features, or debugging facilities.
  \item Executing a reserved encoding raises an illegal-instruction exception.
  \item On an illegal-instruction exception, \code{\%epc} receives the continuation address, with
        the offending address placed into \code{\%edata}.
  \item On an illegal-instruction exception, \code{\%ecause} receives the illegal-instruction exception code.
\end{itemize}

\section{Instruction Set Status}

\begin{itemize}
  \item The VIV-32 v\version{} instruction set is defined by this specification.
        Incompatible future architectural changes require a new architecture version.
  \item Instruction semantics in this chapter are architectural unless explicitly
        described as pseudo-instructions.
  \item Opcode assignments are defined separately by the opcode map.
  \item Assembler syntax and pseudo-instruction expansion rules are defined by the
        toolchain specification.
\end{itemize}
